% Section 6: Conclusion
\section{Conclusion}

\gls{llm} integration into web services represents a shift from parameter-based to intent-based \gls{api}s. This survey examined how natural language interfaces complement RESTful architectures through comparative analysis and AWS Bedrock deployment. Prompt-based \gls{api}s democratize access through natural language but introduce validation, testing, versioning, and security challenges. Non-determinism violates \gls{rest} principles, requiring careful use-case selection. Amazon Nova Micro sufficing for structured extraction at 83\% cost savings versus alternatives \cite{aws_nova_2025} was a critical finding. This principle, matching model capacity to task complexity, is essential for sustainable design.

Security demands layered defences: input sanitization, structural isolation, output guardrails, behavioural constraints \cite{owasp_top_2024, aws_bedrock_2025}. Future architecture will be hybrid: \gls{rest} for deterministic transactions, prompt-based for complex interpretation.

Open challenges: standardization lags adoption, testing methodologies require maturation, token costs create unpredictable expenses, context windows constrain conversational depth \cite{lercher_managing_2024, martinlopez_online_2022}.

For practitioners: start with bounded use cases, leverage managed services, instrument comprehensively, treat prompts as critical artifacts. Decreasing costs and improving capabilities position prompt-based \gls{api}s as complementary technologies enabling natural language interfaces whilst preserving RESTful principles.
